\chapter{Criterio de dossier externo de claims seleccionados}

\section{Dossier externo}
\begin{definition}
Se llama dossier externo de claims seleccionados al paquete disciplinado mediante el cual varios claims ya suficientemente estabilizados, fichables y publicables pueden presentarse fuera del circuito interno como conjunto coherente, trazable y editorialmente sobrio.
\end{definition}

\begin{observation}
El dossier externo no sustituye la tabla madre ni la auditoría interna. Su función es reunir claims ya seleccionados bajo una forma externa más amplia que la ficha mínima, sin caer en apertura indiscriminada del corpus.
\end{observation}

\section{Condición mínima de dossier}
\begin{definition}
Existe dossier externo disciplinado cuando concurren simultáneamente:
\begin{enumerate}
  \item un conjunto de claims ya compatibles con publicación selectiva;
  \item una relación metodológica reconocible entre ellos;
  \item trazabilidad suficiente hacia sus demos o fuentes de origen;
  \item una organización editorial sobria del conjunto;
  \item posibilidad de distinguir, dentro del dossier, entre soporte efectivo, interpretación disciplinada y reservas relevantes.
\end{enumerate}
\end{definition}

\begin{observation}
No se exige todavía una publicación integral del sistema. Se exige solo que el paquete externo no deforme la fuerza real de los claims reunidos.
\end{observation}

\section{Dossier y ficha mínima}
\begin{definition}
Todo dossier externo debe poder descomponerse en fichas públicas mínimas de sus claims constituyentes, sin perder por ello la coherencia del conjunto.
\end{definition}

\begin{observation}
La ficha mínima preserva la unidad; el dossier preserva la relación entre varias unidades.
\end{observation}

\section{Dossier y tabla madre}
\begin{definition}
El dossier externo debe conservar vínculo legible con la tabla madre de demos y claims, de modo que la proyección externa del conjunto no borre la genealogía interna de las piezas reunidas.
\end{definition}

\begin{observation}
Esta cláusula impide que el paquete público funcione como compilación sin raíces metodológicas visibles.
\end{observation}

\section{Condición mínima de no dossier}
\begin{definition}
No debe constituirse dossier externo cuando los claims reunidos solo comparten proximidad retórica, cuando la heterogeneidad del conjunto destruye su sobriedad o cuando el paquete externo induce una lectura más fuerte que la efectivamente alcanzada por sus piezas.
\end{definition}

\begin{observation}
La no constitución de dossier no degrada los claims individuales. Puede significar solo que todavía deben circular por fichas aisladas o permanecer en curación interna.
\end{observation}

\section{Cláusula de alcance}
\begin{observation}
El criterio de dossier externo de claims seleccionados fija una primera forma de agrupación pública sobria dentro del Tomo III. No constituye todavía una política integral de publicación del corpus ni una presentación final exhaustiva del sistema.
\end{observation}

